% Standalone manuscript generated from solution.md with Pandoc.
% Compile twice with XeLaTeX.
\documentclass[10pt,a4paper]{article}
\usepackage[margin=24mm,headheight=15pt]{geometry}
\usepackage{fontspec}
\setmainfont{DejaVuSerif.ttf}[BoldFont=DejaVuSerif-Bold.ttf,ItalicFont=DejaVuSerif-Italic.ttf,BoldItalicFont=DejaVuSerif-BoldItalic.ttf]
\setsansfont{DejaVuSans.ttf}[BoldFont=DejaVuSans-Bold.ttf,ItalicFont=DejaVuSans-Oblique.ttf]
\setmonofont{DejaVuSansMono.ttf}
\usepackage{amsmath,amssymb,unicode-math}
\setmathfont{latinmodern-math.otf}
\usepackage{xcolor,microtype,parskip,needspace,fancyhdr}
\usepackage{bookmark,xurl}
\hypersetup{colorlinks=true,urlcolor=blue,linkcolor=blue,pdftitle={MF-21: the regular expansion threshold for every higher-order zero},pdfauthor={George Stepaniants},pdflang={en-GB}}
\urlstyle{same}
\setlength{\emergencystretch}{3em}
\providecommand{\tightlist}{\setlength{\itemsep}{0pt}\setlength{\parskip}{0pt}}
\providecommand{\pandocbounded}[1]{#1}
\setcounter{secnumdepth}{0}
\allowdisplaybreaks[1]
\widowpenalty=10000
\clubpenalty=10000
\pagestyle{fancy}
\fancyhf{}
\fancyhead[L]{\small\sffamily AFFIRMATIVE RESOLUTION / MF-21}
\fancyhead[R]{\small\sffamily 12 September 2026 (UTC)}
\fancyfoot[L]{\parbox[b]{0.9\textwidth}{\fontsize{8}{10}\selectfont Substantial AI assistance; independent automated review is documented separately.}}
\fancyfoot[R]{\footnotesize\thepage}
\begin{document}
\section*{MF-21: the regular expansion threshold for every higher-order zero}\label{mf-21-the-regular-expansion-threshold-for-every-higher-order-zero}

\textbf{George Stepaniants}\\
Department of Computing and Mathematical Sciences, California Institute of Technology, Pasadena, California, USA\\
12 September 2026 (UTC)

This manuscript proves the complete three-part MF-21 target for every integer \(m\ge3\). The conjecture and its prior special cases are due to Barrera, Böttcher, Grudsky, Maximenko, and the cited later authors. The proof combines an exact boundary determinant with a classical inverse-kernel limit, retaining attribution for that external theorem.

\textbf{Review and assistance.} The \href{https://github.com/ajt60gaibb/OpenProblemsInNLA/blob/main/references/stepaniants-mf21-2026-09-12/independent-review.md}{independent mathematical review} and \href{https://github.com/ajt60gaibb/OpenProblemsInNLA/blob/main/references/stepaniants-mf21-2026-09-12/README.md}{submission record} identify the reviewed source and verification scope. Substantial AI assistance is disclosed. An independent AI-agent audit is informal review; it does not establish external human peer review or formal verification.

\Needspace{10\baselineskip}
\subsection{1. Exact conclusion and the coefficient functions}\label{exact-conclusion-and-the-coefficient-functions}

Fix an integer \(m\ge3\), and put

\nopagebreak[3]
\[
g(\theta)=(2-2\cos\theta)^m,\qquad A_n=T_n(g),\qquad h=(n+2)^{-1},\qquad x_{n,j}=j\pi h.
\]

Let the eigenvalues of \(A_n\), in increasing order and counted with multiplicity, be \(\lambda_{n,j}\).

\textbf{Theorem 1.} There are real \(C^\infty\) functions \(d_0,\ldots,d_{2m}\) on \([0,\pi]\), with \(d_0=g\), such that all three assertions of MF-21 hold. In particular,

\nopagebreak[3]
\[
\left|\lambda_{n,j}-\sum_{k=0}^{p}d_k(x_{n,j})h^k\right|\le C_p h^{p+1}
\quad(0\le p\le2m-1,\ 1\le j\le n)
\tag{1}
\]

for sufficiently large \(n\); the same estimate for \(p=2m\) holds for \(j\ge\lceil(\log(n+2))^2\rceil\); and the latter estimate cannot hold for every \(j\). Constants and lower bounds on \(n\) may depend on \(m,p\), but not on \(n,j\).

Here is a constructive definition of the functions. Write

\nopagebreak[3]
\[
\omega_\ell=e^{2\pi i\ell/m},\qquad 1\le\ell\le m-1.
\]

For \(0<\theta\le\pi\), let \(r_\ell(\theta)\) be the root inside the unit circle of

\nopagebreak[3]
\[
2-r-r^{-1}=\omega_\ell(2-2\cos\theta).
\tag{2}
\]

Define the real phase by

\nopagebreak[3]
\[
\psi(\theta)=\sum_{\ell=1}^{m-1}\arg(1-r_\ell(\theta)e^{-i\theta}),\qquad
\eta(\theta)=\theta+2\psi(\theta),
\tag{3}
\]

where each argument is the principal argument in \((-\pi/2,\pi/2)\). Its extension to \(\theta=0\) will be established below. Extend \(g,\eta\) smoothly to a neighborhood of \([0,\pi]\), and, for sufficiently small real \(h\), let \(Y(x,h)\) be the solution near \(x\) of

\nopagebreak[3]
\[
Y=x+h\eta(Y).
\tag{4}
\]

The coefficients are

\nopagebreak[3]
\[
d_k(x)=\frac{1}{k!}\left.\frac{\partial^k}{\partial h^k}g(Y(x,h))\right|_{h=0}.
\tag{5}
\]

They do not depend on the extension, the matrix size, or the eigenvalue index.

\Needspace{10\baselineskip}
\subsection{2. Stable roots and a normalized exact boundary determinant}\label{stable-roots-and-a-normalized-exact-boundary-determinant}

\textbf{Lemma 2.} The functions \(r_\ell\) extend smoothly to \([0,\pi]\), with \(r_\ell(0)=1\). They are nonzero and satisfy, for some \(c>0\),

\nopagebreak[3]
\[
|r_\ell(\theta)|\le e^{-c\theta}\quad(0\le\theta\le\pi).
\tag{6}
\]

The phase in (3) is smooth and has endpoint values

\nopagebreak[3]
\[
\psi(0)=\frac{(m-1)\pi}{4},\quad \psi(\pi)=0,\quad
\eta(0)=\frac{(m-1)\pi}{2},\quad \eta(\pi)=\pi.
\tag{7}
\]

\emph{Proof.} A root on the unit circle in (2) would make its left side real and nonnegative. This is impossible because \(\omega_\ell\ne1\), including the case \(\omega_\ell=-1\). The reciprocal quadratic has exactly one root inside. Its roots are simple away from \(\theta=0\): a double root could occur only when the right side is 0 or 4.

Set

\nopagebreak[3]
\[
\kappa_\ell=e^{i(\pi\ell/m-\pi/2)}.
\]

The quadratic formula, with its square-root branch chosen near zero, gives

\nopagebreak[3]
\[
r_\ell(\theta)=1-\kappa_\ell\theta+O(\theta^2),\qquad
\Re\kappa_\ell>0.
\tag{8}
\]

The square root has a factor \(\theta\) times an analytic nonvanishing function; this also proves smoothness at zero. The roots never vanish. Compactness, (8), and \(|r_\ell|<1\) away from zero prove (6). All logarithmic derivatives \(r_\ell'/r_\ell\) are bounded.

For positive \(\theta\), every factor in (3) has positive real part. At zero it is \(\theta\) times a smooth nonzero function, with limit

\nopagebreak[3]
\[
\kappa_\ell+i
=2\sin\left(\frac{\pi\ell}{2m}\right)e^{i\pi\ell/(2m)}.
\]

Thus each argument extends smoothly and its limiting value is \(\pi\ell/(2m)\). Summing proves the first value in (7). The multiset of stable roots is closed under conjugation; at \(\theta=\pi\), the factors \(1+r_\ell\) occur in conjugate pairs, with any real factor positive. Their principal arguments sum to zero. This proves the remaining statements. \(\square\)

For \(0<\theta<\pi\), define

\nopagebreak[3]
\[
F_n(\theta)=(n+1)\theta-2\psi(\theta)=(n+2)\theta-\eta(\theta).
\tag{9}
\]

\textbf{Lemma 3.} There are real smooth functions \(E_n\) on \((0,\pi]\) such that \(g(\theta)\) is an eigenvalue of \(A_n\) precisely when

\nopagebreak[3]
\[
\sin F_n(\theta)+E_n(\theta)=0\qquad(0<\theta<\pi).
\tag{10}
\]

For constants \(C,c>0\) depending only on \(m\),

\nopagebreak[3]
\[
|E_n(\theta)|\le C e^{-cn\theta},\qquad
|E_n'(\theta)|\le C(n+1)e^{-cn\theta}.
\tag{11}
\]

Moreover \(E_n(\pi)=0\), and therefore

\nopagebreak[3]
\[
|E_n(\theta)|\le C(n+1)(\pi-\theta)e^{-cn\pi/2}
\qquad(\pi/2\le\theta\le\pi).
\tag{12}
\]

\emph{Proof.} The Laurent polynomial \(g\) is \((2-z-z^{-1})^m\). At the spectral parameter \(g(\theta)\), its \(2m\) distinct characteristic roots are

\nopagebreak[3]
\[
\mathcal Z=(r_1,\ldots,r_{m-1},z,z^{-1},q_1,\ldots,q_{m-1}),\qquad
z=e^{i\theta},\quad q_\ell=r_\ell^{-1}.
\]

An eigenvector is a solution of the corresponding order-\(2m\) recurrence, with ghost values zero at \(1-m,\ldots,0,n+1,\ldots,n+m\). Shifting these indices by \(m-1\), its boundary determinant is

\nopagebreak[3]
\[
D_n(\theta)=\det\begin{pmatrix}
 (w^k)_{0\le k\le m-1,\ w\in\mathcal Z}\\
 (w^{n+m+k})_{0\le k\le m-1,\ w\in\mathcal Z}
\end{pmatrix}.
\tag{13}
\]

The ghost-value criterion is equivalent to the finite eigenvalue equation. Indeed, the characteristic roots form a basis for recurrence solutions. A solution vanishing both in the interior and at the ghost values has at least \(2m\) consecutive zeros and is identically zero. Hence the kernels of the boundary matrix and the eigenvalue matrix have equal dimensions.

For an ordered root list \(U\), put \(V(U)=\prod_{a<b}(U_b-U_a)\). Laplace expansion in the bottom \(m\) rows gives

\nopagebreak[3]
\[
D_n=\sum_{|S|=m}\sigma_S V(S)V(S^c)
             \left(\prod_{w\in S}w\right)^{n+m},\qquad \sigma_S\in\{-1,1\},
\tag{14}
\]

with each sublist in the inherited order. Write \(R=(r_1,\ldots,r_{m-1})\), \(O=(q_1,\ldots,q_{m-1})\),

\nopagebreak[3]
\[
Q=\prod_{\ell=1}^{m-1}q_\ell>0,\qquad
f=\prod_{\ell=1}^{m-1}(1-r_\ell z^{-1})=|f|e^{i\psi}.
\]

Exactly two subsets have maximal product modulus \(Q\): \(S_+=\{z\}\cup O\) and \(S_-=\{z^{-1}\}\cup O\). Their signs are opposite. For some sign \(\sigma\), their coefficient products are

\nopagebreak[3]
\[
\sigma V(R)V(O)Q z^{-(m-1)}\overline f^{2},\qquad
-\sigma V(R)V(O)Q z^{m-1}f^2.
\tag{15}
\]

For example, this follows by multiplying \(V(z,O)=V(O)\prod(q_\ell-z)\) and \(V(R,z^{-1})=V(R)\prod(z^{-1}-r_\ell)\). Their contribution to (14) is consequently

\nopagebreak[3]
\[
\mathcal N_n(\theta)\sin F_n(\theta),\qquad
\mathcal N_n=2i\sigma V(R)V(O)Q^{n+m+1}|f|^2.
\tag{16}
\]

This normalizer is nonzero for \(0<\theta\le\pi\). Every other subset has

\nopagebreak[3]
\[
\left|b_S(\theta)\right|:=\left|\frac{\prod_{w\in S}w}{Q}\right|
\le \max_\ell|r_\ell(\theta)|\le e^{-c\theta}.
\tag{17}
\]

To see the first inequality, if a subset has both oscillatory roots, it omits at least one exterior root; if it has neither, it contains at least one interior root; and if it has exactly one, it replaces at least one exterior root by an interior root. Each replacement only decreases the modulus further.

Divide every remaining term of (14) by (16). This expresses \(E_n\) as a finite sum

\nopagebreak[3]
\[
E_n(\theta)=\sum_{S\ne S_+,S_-}a_S(\theta)b_S(\theta)^{n+m},
\qquad
a_S=\frac{\sigma_S V(S)V(S^c)}{2i\sigma V(R)V(O)Q|f|^2}.
\tag{18}
\]

All \(a_S\) and their first derivatives are bounded on \([0,\pi]\). At zero, every pairwise root difference has a simple zero: the \(2m\) first derivatives of the roots are the distinct numbers \(-\kappa_\ell,\kappa_\ell,i,-i\). Thus both the numerator and denominator of \(a_S\) vanish to exactly order \(m(m-1)\), with nonzero leading coefficients. Their quotient extends smoothly. At \(\pi\), the denominator has no zero: the interior roots and the exterior roots are separately distinct, and \(f\ne0\). Some numerators vanish there because \(z=z^{-1}\), which causes no singularity.

Every \(b_S\) is smooth, nonzero, and has bounded logarithmic derivative. Differentiating (18) and using (17) proves (11), after changing constants. The quotient \(D_n/\mathcal N_n\) is real for real \(\theta\): conjugation permutes the two oscillatory roots, and permutes the interior and exterior lists by the same permutation. Hence \(\overline D_n=-D_n\), while \(V(R)V(O)\), \(Q\), and \(|f|^2\) are real, so \(\overline{\mathcal N_n}=-\mathcal N_n\). This proves that \(E_n\) is real.

At \(\pi\), the two oscillatory columns in (13) coincide, so \(D_n(\pi)=0\), whereas \(\mathcal N_n(\pi)\ne0\) and \(\sin F_n(\pi)=\sin((n+1)\pi)=0\). Therefore \(E_n(\pi)=0\). Integrating the derivative bound over \([\theta,\pi]\) proves (12). \(\square\)

\Needspace{10\baselineskip}
\subsection{3. Root indexing and a uniform exponentially small phase error}\label{root-indexing-and-a-uniform-exponentially-small-phase-error}

All eigenvalues of \(A_n\) lie strictly between 0 and \(4^m\). For nonzero \(v\), its quadratic form is the integral of \(g\) times the squared modulus of a nonzero trigonometric polynomial; both \(g\) and \(4^m-g\) are positive almost everywhere. Hence there are exactly \(n\) eigenangles in \((0,\pi)\), counted with multiplicity.

\textbf{Lemma 4.} There is an integer \(J\ge1\), depending only on \(m\), such that for all sufficiently large \(n\) and all \(J\le j\le n\), the \(j\)-th eigenangle \(\theta_{n,j}=g^{-1}(\lambda_{n,j})\) is simple and satisfies

\nopagebreak[3]
\[
|\theta_{n,j}-Y(x_{n,j},h)|\le \frac{C e^{-cj}}{n+2},\qquad
0<\theta_{n,j},Y(x_{n,j},h)\le C\frac{j}{n+2}.
\tag{19}
\]

\emph{Proof.} Boundedness of \(\eta'\) gives, for large \(n\),

\nopagebreak[3]
\[
\frac{n+2}{2}\le F_n'(\theta)\le2(n+2),\qquad
F_n(0)=-\frac{(m-1)\pi}{2},\quad F_n(\pi)=(n+1)\pi.
\tag{20}
\]

Let \(I_{n,k}\) be the interval on which \(|F_n(\theta)-k\pi|\le\pi/4\), for \(J\le k\le n\). Choose \(J\) large enough. On the whole range \(F_n\ge J\pi-\pi/4\), boundedness of \(\eta\) implies \(n\theta\ge c_1J-C_1\). Thus (11) makes \(|E_n|<1/4\) and \(|E_n'|<(n+2)/8\) throughout this range, uniformly for all sufficiently large \(n\).

At the two endpoints of \(I_{n,k}\), the sine has opposite signs and magnitude \(1/\sqrt2\). On \(I_{n,k}\), its derivative has constant sign and magnitude at least \((n+2)/(2\sqrt2)\), so the derivative of (10)'s left side has that same sign. There is exactly one root in each \(I_{n,k}\), and it is a simple zero of the determinant. The corresponding eigenspace has dimension one: otherwise the boundary matrix in (13) would have nullity at least two, making its determinant derivative zero. Since \(A_n\) is symmetric, this eigenvalue also has algebraic multiplicity one.

Between these intervals, and below the final interval \(|F_n-(n+1)\pi|\le\pi/4\), the magnitude of the sine is at least \(1/\sqrt2\); there are no roots there. The final interval needs a different argument because \(\theta=\pi\) is an artificial determinant root. There, \(\theta\ge\pi/2\) for large \(n\), and (20) and the elementary lower bound for sine give

\nopagebreak[3]
\[
|\sin F_n(\theta)|\ge c_2(n+2)(\pi-\theta).
\]

Equation (12) is strictly smaller for every \(\theta<\pi\), once \(n\) is large. Thus there is no eigenangle in this final interval.

We have found exactly one simple eigenangle in each interval \(I_{n,k}\), \(k=J,\ldots,n\), and none between or above them. Counting downwards from the largest eigenvalue therefore identifies the root in \(I_{n,k}\) as precisely the \(k\)-th eigenangle; no assertion about the roots below this range was needed.

At that root, (11) and \(|F_n-j\pi|\le\pi/4\) imply

\nopagebreak[3]
\[
|F_n(\theta_{n,j})-j\pi|\le C e^{-cn\theta_{n,j}}\le C e^{-c'j}.
\]

Here \(n\theta_{n,j}\ge c_3j-C_3\), by boundedness of \(\eta\). The exact root \(F_n(y)=j\pi\) is \(y=Y(x_{n,j},h)\). Applying (20) proves the first bound in (19). The upper bounds follow from \(F_n=(n+2)\theta-\eta(\theta)\), boundedness of \(\eta\), and \(j\ge J\). \(\square\)

We also require elementary estimates covering the few remaining eigenvalues. Let \(N=n+2m\), and let \(C_N\) be the circulant matrix with eigenvalues \(g(2\pi\ell/N)\), \(0\le\ell<N\). Its leading \(n\times n\) principal submatrix is \(A_n\), because no Fourier coefficient of bandwidth \(m\) wraps into that block. Its increasingly ordered eigenvalues are

\nopagebreak[3]
\[
\mu_{N,j}=g\left(\frac{2\pi\lfloor j/2\rfloor}{N}\right),\qquad 1\le j\le N.
\]

Cauchy interlacing gives

\nopagebreak[3]
\[
\mu_{N,j}\le\lambda_{n,j}\le\mu_{N,j+2m}.
\tag{21}
\]

Consequently \(\lambda_{n,j}\le C_m(j+2m)^{2m}n^{-2m}\). Also, for \(j\ge2\), using \(2\sin(t/2)\ge2t/\pi\) on \([0,\pi]\),

\nopagebreak[3]
\[
\lambda_{n,j}\ge\left(\frac{4\lfloor j/2\rfloor}{n+2m}\right)^{2m}
\ge\left(\frac{4j}{3(n+2m)}\right)^{2m}.
\tag{22}
\]

\Needspace{10\baselineskip}
\subsection{4. The expansions through the claimed threshold}\label{the-expansions-through-the-claimed-threshold}

The implicit function theorem applied uniformly on the compact interval in (4) gives smooth \(Y\), so for any fixed \(p\),

\nopagebreak[3]
\[
g(Y(x,h))=\sum_{k=0}^p d_k(x)h^k+O(h^{p+1})
\tag{23}
\]

uniformly in \(x\in[0,\pi]\). Since \(g\) vanishes to order \(2m\) at zero, repeated differentiation in \(h\) shows

\nopagebreak[3]
\[
d_k(x)=O(x^{2m-k})\quad(0\le k\le2m),\qquad d_0=g.
\tag{24}
\]

Indeed, every term in the \(k\)-th derivative is a bounded smooth factor times \(g^{(s)}(x)\) for \(s\le k\), and \(g^{(s)}(x)=O(x^{2m-s})\). For \(k=0\) the same conclusion is immediate.

For \(j\ge J\), the mean value theorem, \(g'(\theta)=O(\theta^{2m-1})\), and (19) yield

\nopagebreak[3]
\[
|\lambda_{n,j}-g(Y(x_{n,j},h))|
\le C h^{2m}j^{2m-1}e^{-cj}\le C h^{2m}.
\tag{25}
\]

For \(j<J\), (21) gives \(\lambda_{n,j}=O(h^{2m})\), and (24) gives \(\sum_{k=0}^{2m-1}d_k(x_{n,j})h^k=O(h^{2m})\), uniformly over this fixed finite set. Combining these observations with (23) at \(p=2m-1\) proves (1) at that order for every \(j\). Subtracting the omitted terms, whose coefficients are bounded on \([0,\pi]\), proves every lower order in (1), with the same \(d_k\).

If \(j\ge\lceil(\log(n+2))^2\rceil\), the first bound in (25) is \(O(h^{2m+1})\): exponential decay beats the displayed polynomial uniformly in that range. Equation (23) with \(p=2m\) now proves Part 2 of the target.

\Needspace{10\baselineskip}
\subsection{\texorpdfstring{5. A trace obstruction to extending order \(2m\) to the lowest indices}{5. A trace obstruction to extending order 2m to the lowest indices}}\label{a-trace-obstruction-to-extending-order-2m-to-the-lowest-indices}

We use a classical inverse-kernel limit, rather than an unproved formula for each extreme eigenvalue. For \(A_n=T_n(|1-z|^{2m})\), the known limit is

\nopagebreak[3]
\[
n^{1-2m}(A_n^{-1})_{\lceil nx\rceil,\lceil ny\rceil}\longrightarrow G_m(x,y)
\quad\hbox{in }L^\infty([0,1]^2),
\tag{26}
\]

with endpoint indices interpreted in \(\{1,\ldots,n\}\). The continuous Green kernel is that of \((-1)^m d^{2m}/dx^{2m}\) with derivatives of orders \(0,\ldots,m-1\) vanishing at both endpoints. The precise input (26), including the explicit kernel, is recorded in Böttcher--Widom {[}BW{]}, §2, pp.~3--4, in the paragraph preceding (13); its formula (5) gives, when \(x+y\ge1\),

\nopagebreak[3]
\[
G_m(x,y)=\frac{x^m y^m}{((m-1)!)^2}
\int_{\max(x,y)}^1\frac{(t-x)^{m-1}(t-y)^{m-1}}{t^{2m}}\,dt.
\tag{27}
\]

The kernel is invariant under \((x,y)\mapsto(1-x,1-y)\). Substituting \(u=1-x/t\) in (27) at \(y=x\ge1/2\), and then using this symmetry, gives on all of \([0,1]\)

\nopagebreak[3]
\[
G_m(x,x)=\frac{x^{2m-1}(1-x)^{2m-1}}{(2m-1)((m-1)!)^2}.
\tag{28}
\]

The essential-uniform convergence in (26) also controls the diagonal: the approximating kernel is constant on each grid square, and continuity of \(G_m\) bounds the difference between an interior point of such a square and its diagonal. Thus integrating the diagonal yields

\nopagebreak[3]
\[
\lim_{n\to\infty}(n+2)^{-2m}\operatorname{tr}(A_n^{-1})
=\int_0^1G_m(x,x)\,dx
=\frac{((2m-1)!)^2}{(4m-1)!(2m-1)((m-1)!)^2}
\in\mathbb Q.
\tag{29}
\]

Suppose, contrary to Part 3, that the estimate of order \(2m\) were uniform for every \(j\). For each fixed \(j\), (23) then gives

\nopagebreak[3]
\[
\lambda_{n,j}=g(Y(\pi jh,h))+O(h^{2m+1}).
\]

From (4) and (7), \(Y(\pi jh,h)/h\to\pi j+(m-1)\pi/2\). Hence, writing \(a=(m-1)/2\),

\nopagebreak[3]
\[
h^{-2m}\lambda_{n,j}\longrightarrow \pi^{2m}(j+a)^{2m}.
\tag{30}
\]

The inverse terms \(h^{2m}/\lambda_{n,j}\), extended as zero for \(j>n\), are bounded by \(C_m j^{-2m}\) for \(j\ge2\), by (22). The \(j=1\) term is eventually bounded by (30). Dominated convergence for the counting measure therefore gives a second value for (29):

\nopagebreak[3]
\[
\lim_{n\to\infty}h^{2m}\operatorname{tr}(A_n^{-1})
=\pi^{-2m}\sum_{j=1}^{\infty}(j+a)^{-2m}.
\tag{31}
\]

This number is not rational. If \(m=2r+1\ge3\), the sum before multiplication by \(\pi^{-2m}\) is

\nopagebreak[3]
\[
\zeta(2m)-\sum_{\ell=1}^{r}\ell^{-2m}.
\tag{32}
\]

If \(m=2r\ge4\), it is

\nopagebreak[3]
\[
(2^{2m}-1)\zeta(2m)-\sum_{\ell=0}^{r-1}(\ell+1/2)^{-2m}.
\tag{33}
\]

Euler's formula makes \(\zeta(2m)/\pi^{2m}\) rational. In both (32) and (33), the subtracted finite sum is a nonzero positive rational number. Thus (31) has the form \(u-v/\pi^{2m}\), where \(u,v\in\mathbb Q\) and \(v>0\). Transcendence of \(\pi\) implies this is irrational. This contradicts (29), proving Part 3 and completing Theorem 1. \(\square\)

The contradiction does not rely on an interchange of a merely pointwise spectral limit without domination; the uniform tail majorant was supplied explicitly by (22). It also does not infer a general-\(m\) statement from the \(m=2\) or \(m=3\) cases.

\Needspace{10\baselineskip}
\subsection{6. Sources, attribution, and verification scope}\label{sources-attribution-and-verification-scope}

The target is Conjecture 8.4 of Barrera, Böttcher, Grudsky, and Maximenko {[}BBGM{]}, p.~26. Their Theorem 1.2 establishes the analogous \(m=2\) threshold. The later seven-diagonal work {[}BGSV{]} treats \(m=3\); the 2026 survey {[}B26{]}, section ``Beyond the simple-loop class,'' distinguishes these results from broader local expansions. These contributions retain their original attribution.

The only substantive external Toeplitz theorem used above is the established uniform inverse-kernel limit (26), with the kernel formula (27). The determinant argument and its error estimates, endpoint counting, all-order Taylor construction, and trace contradiction are proved here. Cauchy interlacing, the finite-dimensional implicit function theorem, Euler's even-zeta formula, and transcendence of \(\pi\) are standard mathematical inputs. No numerical calculation is a premise of the proof.

The linked independent review and public eligibility audit record the verification status and bounded search scope. The source search checked the original conjecture, the seven-diagonal preprint, the 2026 survey, and targeted later work; it found no full resolution of the three parts for every \(m\ge3\). A bounded public search does not certify novelty or exclude private or unpublished work. No formal verification, external human peer review, or priority certification is asserted.

\textbf{{[}BBGM{]}} M. Barrera, A. Böttcher, S. M. Grudsky, and E. A. Maximenko, \emph{Eigenvalues of even very nice Toeplitz matrices can be unexpectedly erratic}, Operator Theory: Advances and Applications 268 (2018), 51--77. \href{https://arxiv.org/abs/1710.05243}{arXiv:1710.05243}, Conjecture 8.4 and Theorem 1.2; \href{https://doi.org/10.1007/978-3-319-75996-8_2}{DOI}.

\textbf{{[}BW{]}} A. Böttcher and H. Widom, \emph{From Toeplitz eigenvalues through Green's kernels to higher-order Wirtinger--Sobolev inequalities}, Operator Theory: Advances and Applications 171 (2006), 73--87. \href{https://arxiv.org/abs/math/0412269}{arXiv:math/0412269}, formula (5), §2 paragraph preceding (13), and (13); \href{https://doi.org/10.1007/978-3-7643-7980-3_4}{DOI}.

\textbf{{[}BGSV{]}} M. Barrera, S. Grudsky, V. Stukopin, and I. Voronin, \emph{Asymptotics of the eigenvalues of seven-diagonal Toeplitz matrices of a special form}, Advances in Operator Theory 9 (2024), 79. \href{https://arxiv.org/abs/2111.07196}{arXiv:2111.07196}, Theorems 2.3--2.6; \href{https://doi.org/10.1007/s43036-024-00374-1}{DOI}.

\textbf{{[}B26{]}} A. Böttcher, \emph{Ten years with Sergei Grudsky in the eigenvalue bulk of Toeplitz matrices}, Journal of Mathematical Sciences 298 (2026), 363--376. \href{https://doi.org/10.1007/s10958-025-07833-x}{DOI}, section ``Beyond the simple-loop class.''

\end{document}
